\section{Conclusions}\label{conc}
\subsection{Project summary}
\subsection{\nmu Project \nmu limitations}\label{conc:project_limits}
In an ideal world, it would have been possible to have created a project that could target any and all \latexspaced accessibility concerns. Unfortunately, that is outwith the scope of the project, which has been defined as follows:
\begin{itemize}
    \item \lintexspaced is not designed to handle the majority of parsing. Instead, it focuses on a few key elements that are related to the accessibility concerns targetted in our context. Otherwise, it handles a small amount of verification regarding the structure of individual lines and the overall document structure.
    \item This project does not aim to provide automated fixes for errors encountered.
    \item \lintexspaced does not contain an in-built screen reader. While some considerations were put aside for how this could be accomplished, it is not feasible to create an in-built screen reader with the timeframe allocated for this project, due to the vast considerations needed. It might even be a fair statement to consider it a project on its own.
    \item \lintexspaced is not intended to find usage of outdated or incorrect accessibility, but rather to enhance a document's screen reader accessibility through additions to the document.
    \item \lintexspaced is not intended to debug \latexspaced code, as a dedicated \acrfull{ide} and \texspaced engine should be used for this instead.
    \item \lintexspaced expects code to compile in \lualatexspaced before its changes have been applied, with other \texspaced engines being considered outdated and consequently not supported for this context.
    \item \lintexspaced places emphasis on accessibility features for screen reader users.
    \item It is unfeasible to account for all \gls{tagging} options in our context. Therefore, a few key options have been selected based on what would be expected to have the most impact.
\end{itemize}

\subsection{Achievements}


\subsection{Future work}
\lintexspaced has the potential to grow in a variety of manners, such as targetting the limitations in \fullref{conc:project_limits}.

An in-built screen reader for \lintexspaced is an interesting proposition, with it likely being a project by itself due to the lack of finer details requiring vast amounts of consideration and the difficulties with finding similar materials, seeing as most references for an in-built screen reader appear to point to screen readers as a browser extension. However, it may be possible to make \lintexspaced work with those extensions to provide the ideal experience for the end user.

Automated fixes, with the addition of being able to fix several identical errors at once, would be a great feature to include in \lintex . There is potential for this feature to work alongside the error handling of \lintex , yet this has unfortunately not been explored further at this time.

Additionally, highlighting why \lintexspaced has recommended a change would make the experience a lot more insightful and better encourage developers to be inclusive of screen reader users in the future.

Finally, utilising HTMX to make \lintexspaced more responsive to the end user's actions would allow for a more immersive experience for the user, while keeping the technologies required at a reasonable minimum requirements.
\subsection{Reflection}


\subsection{Final thoughts}